\chapter*{Prolog.}
\section*{The Tome Eternal Struggle.}

And the Lord spake, saying:

“Behold, I giveth unto thee the Scrolls of Understanding, summoned forth from the depths of \LaTeX\ by the hand of My servant. Yea, these are no mere scribbles — but a mighty Halberd to smite the Demon of Academic Failure, a Rod of Iron against the Serpent of Procrastination, and a Lamp unto thy feet in the dark valley of the Night Before the Exam, when gnashing of teeth shall be heard throughout the land.”

“These Scrolls were forged for this very battle, and for no other.”

And lo, the Lord continued:

“Verily, if thou shalt read these Scrolls with diligence and trembling, the veils shall be lifted: the Sacred Structures of Group, Ring, and Field shall be revealed in glory. No longer shall homomorphisms confound thee, nor shall normal subgroups flee thy grasp. The Lamentations of those who knew not the Sylow Theorems shall not be thine.”

“Even the Lemma of Zorn, whose name is spoken in fear, shall not consume thee. For thou shalt wield the Axiom of Choice as a sword, and thy understanding shall be complete.”

“Yea, though thou walkest through the Valley of Final Examinations, thou shalt fear no theorem, for thy Scrolls are with thee; thy pen and thy highlighter, they comfort thee.”

“Go forth now, child of Algebra, and multiply... within a ring. And may thine direct product be ever internal.”

\section*{Here Am I. Send me.}
\textit{And I heard the voice of the Lord, saying,
“Whom shall I send, and who will go for us?”
Then said I, “Here am I; send me.”} 

\section*{Words of truth.}
Tylko pamiętaj, mów ten monolog tak, jak ja ci go przeczytałem: płynnie i swobodnie. Gdybyś miał deklamować z fałszywym patosem, jak niektórzy aktorzy, wolałbym już, żeby moje wiersze wyryczał herold miejski. I nie machaj rękami, jakbyś drwa rąbał. Używaj gestów oszczędnie: w najdzikszym zamęcie, wirze, nawałnicy namiętności trzeba dyscypliny - bez niej nie da się pasji składnie wyrazić.

